\Legend
记字母表大小为 $sigma=5$。对于一个询问，开头字符 $S_l$ 和末尾字符 $S_r$ 已经固定。中间只能从 $S_{l+1}ldots S_{r-1}$ 取子序列，且允许取空串。每个不同的中间字符串恰好对应一个合法答案。

\textbf{按字典序选择}

在已经选择了某个前缀、下一次可以从位置 $p$ 开始选时，设 $f_c(p)$ 是区间 $S_pldots S_{r-1}$ 中以字符 $c$ 开头的不同非空子序列数。给每个字符串选最靠左的下标表示，则所有首字符为 $c$ 的字符串一一对应于选定第一次出现的 $c$ 后面的不同子序列。因此，如果 $j$ 是不小于 $p$ 的第一个 $c$，有

$$f_c(p)=F(j+1),\qquad F(p)=1+\sum_c f_c(p),$$

其中 $F(p)$ 计入空子序列；不存在这样的 $j$ 时 $f_c(p)=0$。

下一字符为 $c$ 的答案共有 $f_c(p)+[c=S_r]$ 个。额外的 $1$ 是直接把固定末尾字符接到当前前缀之后，形成的较短字符串。按 $c$ 从小到大扣除各分支的数量，选出包含第 $k$ 个答案的分支；当 $c=S_r$ 时，直接结束的那个字符串排在该分支的最前面。任何数量只需保存到 $10^9+1$，因为 $k\le10^9$。

\textbf{区间计数矩阵}

用六维向量 $(f_a,f_b,f_c,f_d,f_e,1)$ 表示后缀状态。从右往左读到字符 $x$ 时，只有 $f_x$ 改为 $1+f_a+f_b+f_c+f_d+f_e$，其余分量不变。这是一个系数非负的 $6\times6$ 线性变换。

在线段树的每个结点保存对应区间的变换矩阵。从右边界的 $(0,0,0,0,0,1)$ 出发，按右到左顺序应用覆盖 $[p,r)$ 的 $O(\log N)$ 个结点，便能求出全部 $f_c(p)$。矩阵和向量的运算全部在 $10^9+1$ 截断；因为系数非负，截断不会改变与 $k$ 的比较结果。

\textbf{一次跳过同字母段}

逐字符选答案可能经过 $10^5$ 个相同字母。设接下来重复选择字符 $c$，按位置依次记这些 $c$ 为 $c_1,c_2,\ldots$。在选第 $j$ 个 $c$ 前，所有较小字符的分支数之和为 $L_j$。还须跳过直接以末尾字符结束的短字符串；因此每选择一个 $c$，从当前排名中扣去

$$L_j+[S_r\le c].$$

若已经选择前 $t$ 个 $c$，累计扣数为 $D_t$，最后一个 $c$ 后面的不同子序列数（包括空串）为 $F_t$，那么能继续选择这 $t$ 个字符当且仅当

$$D_t<k\le D_t+F_t.$$

左界随 $t$ 不减，右界随 $t$ 不增：从上一步的上界中扣掉所有大于 $c$ 的分支便得到下一步的上界。所以可在线段树上从左向右跳过整块 $c$，只在第一块不满足条件的结点内部下降。

结点除原矩阵外，为每个 $c$ 再保存几个六维行向量：区间内每次经过 $c$ 后较小字符分支数的总和、去掉最后一次的总和、最后一个 $c$ 后的较小字符分支数，以及最后一个 $c$ 后的 $F$。它们都能由左右儿子的行向量与区间矩阵合并。检查一个完整结点时，利用去掉最后一次的总和与进入该结点前的较小分支数，就能得到新的 $D_t$；利用最后一个 $c$ 后的 $F$ 检查上式。每次跳段需 $O(\log N)$ 个结点、每个结点常数次六维运算。哈希值用预处理的 $1331^t$ 和 $1+1331+\cdots+1331^{t-1}$ 一次加入整段。

\textbf{段数与复杂度}

设所选中间字符串有 $T$ 个同字母段，其中下降的相邻段边界有 $D$ 个。从这些下降边界中选奇数位或偶数位的一组，至少有 $\lceil D/2\rceil$ 个且互不相邻。任意删除其中若干个下降边界之前的整段，仍得到原区间的合法子序列。对两个不同的删除集合，在第一个不同选择的整段处，一个字符串保留该段，另一个立即出现更小的下一段字符，因此得到不同且更小的字符串。故 $k\ge2^{\lceil D/2\rceil}$，即 $D\le2\log_2 k$。

字符只有五种。相邻段严格上升的次数至多为 $4D+4$，所以 $T\le5D+5=O(\sigma\log k)$。因此逐询问至多进行 $O(\sigma\log k)$ 次跳段。

线段树预处理时间 $O(N\sigma^3)$、空间 $O(N\sigma^2)$；单次询问时间 $O(\sigma^3\log k\log N)$。这里 $\sigma=5$ 是固定常数。主解使用无符号 $32$ 位整数自然溢出实现模 $2^{32}$ 哈希。

\tt{solution2.cpp} 和 \tt{solution.py} 使用逐询问直接构建后缀计数的独立算法，只在小数据组中参与交叉验证。
